%%%%%%%%%%%%%%%%%
% CV targeted for Arm - System Performance Analysis Engineer
%%%%%%%%%%%%%%%%%

\documentclass[8pt,a4paper,ragged2e,withhyper]{altacv}

\geometry{left=1.25cm,right=1.25cm,top=.5cm,bottom=1.cm,columnsep=1.2cm}

\usepackage{paracol}
\usepackage{fontawesome5}
\usepackage{hyperref}

\iftutex 
  \setmainfont{Roboto Slab}
  \setsansfont{Lato}
  \renewcommand{\familydefault}{\sfdefault}
\else
  \usepackage[rm]{roboto}
  \usepackage[defaultsans]{lato}
  \renewcommand{\familydefault}{\sfdefault}
\fi

\definecolor{SlateGrey}{HTML}{2E2E2E}
\definecolor{LightGrey}{HTML}{666666}
\definecolor{DarkPastelRed}{HTML}{8AC0D9}
\definecolor{PastelRed}{HTML}{00B0DA}
\definecolor{GoldenEarth}{HTML}{B4AD0C}

\colorlet{name}{black}
\colorlet{tagline}{PastelRed}
\colorlet{heading}{DarkPastelRed}
\colorlet{headingrule}{GoldenEarth}
\colorlet{subheading}{PastelRed}
\colorlet{accent}{PastelRed}
\colorlet{emphasis}{SlateGrey}
\colorlet{body}{LightGrey}

\definecolor{violet}{HTML}{5A2582}
\definecolor{rouge}{HTML}{F53B3B}
\definecolor{noir}{HTML}{00232C}
\definecolor{bleu}{HTML}{00B0DA}
\definecolor{rose}{HTML}{F831A0}
\definecolor{jaune}{HTML}{FFEC00}
\definecolor{vert}{HTML}{34AD6C}

\renewcommand{\namefont}{\Huge\rmfamily\bfseries}
\renewcommand{\personalinfofont}{\footnotesize}
\renewcommand{\cvsectionfont}{\LARGE\rmfamily\bfseries}
\renewcommand{\cvsubsectionfont}{\large\bfseries}

\renewcommand{\cvItemMarker}{{\small\textbullet}}
\renewcommand{\cvRatingMarker}{\faCircle}

\input{../../_shared/pubs-num.tex}
\addbibresource{../../_shared/sample.bib}

\begin{document}

\name{BOILEAU Guillaume}
\tagline{System Performance Analysis Engineer | Modelling, Statistics, Diagnostics \& Complex Systems}

\photoL{2.8cm}{../../_shared/moi}

\personalinfo{%
  \email{guillaume.boileaupro@gmail.com}
  \phone{+33 6 59 94 85 84}
  \location{Alpes-Maritimes, FRANCE}
  \linkedin{guillaume-boileau-phd-891b81265}
  \researchgate{Guillaume-Boileau-2}
  \orcid{0000-0002-3576-6968}
}

\makecvheader

\vspace{-0.3cm}
\AtBeginEnvironment{itemize}{\small}

\columnratio{0.64}

\begin{paracol}{2}

\cvsection{Experience}

\cvevent{\textbf{Software Engineer -- System Integration, Operational Performance \& Production Diagnostics}}{VEV Platform Services France}{2025 -- 2026}{Valbonne, France}

\begin{itemize}
\item Developed and maintained web and backend services for an EV charging CPMS platform using \textbf{TypeScript, Angular and Node.js}.
\item Contributed to software integration activities involving operational charging infrastructure, hardware interfaces and field equipment constraints.
\item Analysed software behaviours, operational data exchanges, communication issues and service anomalies affecting production continuity.
\item Investigated data inconsistencies, abnormal behaviours and system-level interactions between software services and connected charging equipment.
\item Supported monitoring, verification, correction and follow-up of operational data flows, including FTP-based exchanges.
\item Contributed to issue diagnosis, prioritization and technical follow-up with stakeholders in an industrial operational context.
\item Strengthened experience with production constraints, maintainability, service robustness and software systems connected to physical infrastructure.
\end{itemize}

\divider

\cvevent{\textbf{Senior R\&D Engineer -- System Performance, Modelling, Validation \& Technical Coordination}}{Laboratoire Lagrange, Observatoire de la Côte d'Azur, CNES}{2023 -- 2025}{Nice, France}

\begin{itemize}
\item Contributed to engineering and analysis activities for the \textbf{LISA ESA/NASA space mission}, involving complex system-level performance constraints, scientific requirements and high modelling rigor.
\item Developed \textcolor{DarkPastelRed}{\href{https://gitlab.oca.eu/gboileau/lisa_l3_tools}{\bf CATpyLISA}}, a Python platform for large-scale model integration, comparison, analysis and validation workflows.
\item Conducted deep technical studies on complex system behaviours, weak-signal environments, numerical models, uncertainty sources and performance limitations.
\item Analysed sensitivity to modelling assumptions, numerical parameters, data-processing choices, configuration changes and validation conditions.
\item Identified bottlenecks, inconsistencies, deviations, critical behaviours and limiting factors in complex software and analysis chains.
\item Structured validation workflows to compare outputs, verify expected behaviours, detect deviations and consolidate technical evidence.
\item Supported activities aligned with \textbf{V-cycle engineering logic}: requirements interpretation, validation planning, consistency checks, result verification and anomaly analysis.
\item Produced technical analyses supporting model selection, risk identification, performance interpretation and decision-making.
\item Maintained traceability of configurations, assumptions, parameters, validation results and technical conclusions.
\item Coordinated technical activities involving contributors, research engineers, technicians and Master's thesis interns.
\item Operated in a collaborative international framework requiring technical English, structured reporting, versioned workflows and rigorous documentation.
\end{itemize}

\divider

\cvevent{\textbf{Data Scientist -- Instrument Performance, Noise Modelling \& Statistical Diagnostics}}{Universiteit Antwerpen, Belgium}{2021 -- 2023}{Antwerp, Belgium}

\begin{itemize}
\item Analysed instrumental and environmental effects impacting the sensitivity, stability and performance of precision measurement systems.
\item Developed Python-based analysis pipelines for noise source identification, uncertainty estimation, Bayesian inference and statistical modelling.
\item Conducted deep technical studies on degradation patterns, dominant contributors, system limitations and correlations between environmental conditions and measurement performance.
\item Investigated unexpected behaviours, weak points and possible mitigation directions using statistical diagnostics and reproducible workflows.
\item Built traceable analysis workflows using Unix environments, scripting, source control practices and structured documentation.
\item Contributed to studies connected to gravitational-wave detector systems, where performance, noise control, robustness and system-level sensitivity are critical.
\item Communicated complex technical results to scientific and engineering stakeholders through reporting, presentations and collaborative review.
\end{itemize}

\divider
\newpage

\cvevent{\textbf{Junior R\&D Engineer -- Signal Analysis, Simulation \& System Performance Assessment}}{ARTEMIS, Observatoire de la Côte d'Azur}{2018 -- 2021}{Nice, France}

\begin{itemize}
\item Conducted research and engineering-oriented analysis for the preparation of gravitational-wave space missions, with a strong focus on \textbf{LISA}.
\item Developed methods for the separation of weak overlapping signals in complex low-SNR datasets.
\item Built simulation and analysis workflows to evaluate system performance under realistic noise, foreground and uncertainty conditions.
\item Compared simulated outputs, model behaviours and analysis results to assess consistency, robustness and performance limitations.
\item Investigated noise sources, correlations and sensitivity limits through diagnostic analyses relevant to precision instruments.
\item Produced peer-reviewed results contributing to the understanding of performance limits for future gravitational-wave space missions.
\item Developed strong expertise in modelling, uncertainty analysis, signal processing, technical validation and scientific software workflows.
\end{itemize}

\cvsection{Professional Training}

\cvevent{Supervised Machine Learning}{DeepLearning.AI -- Andrew Ng}{2020}{Coursera}

\divider

\cvevent{Recherche reproductible : principes méthodologiques pour une science transparente}{FUN MOOC -- Inria}{2025}{France Université Numérique}

\begin{itemize}
\item Training focused on reproducible research practices, computational documentation, data management, version control and transparent analysis workflows.
\item Practical exposure to \textbf{Git/GitLab}, \textbf{Jupyter notebooks}, structured documentation, data handling and reproducible software environments.
\end{itemize}

\divider

\cvevent{Continuous Professional Training -- Software, Data, AI and Systems}{OpenClassrooms}{2024 -- 2026}{On line}

\begin{itemize}
\item Completed courses covering \textbf{Python}, \textbf{JavaScript}, \textbf{TypeScript}, \textbf{Angular}, \textbf{Node.js}, \textbf{Django}, \textbf{REST APIs}, \textbf{SQL}, \textbf{Machine Learning}, \textbf{AI}, \textbf{Linux}, \textbf{C/C++}, web development and software engineering fundamentals.
\item Strengthened practical skills in full-stack development, data analysis, algorithmic thinking, database querying, Linux administration and modern software tooling.
\end{itemize}

\switchcolumn

\cvsection{Summary and Strengths}

\textbf{System performance and technical analysis engineer with strong experience in modelling, statistical diagnostics, simulation workflows, software/data validation and complex system analysis}. Background developed in international space and precision-instrumentation environments, complemented by recent industrial software experience involving operational data flows, hardware/software interfaces and production diagnostics.

Experienced in \textbf{deep technical studies}, \textbf{performance analysis}, \textbf{statistical modelling}, \textbf{uncertainty analysis}, \textbf{bottleneck identification}, \textbf{anomaly investigation}, \textbf{Python automation}, \textbf{Unix/Linux environments}, \textbf{scripting}, \textbf{source control}, \textbf{technical documentation} and \textbf{cross-functional coordination}.

Able to analyse complex systems, identify limiting factors, consolidate technical evidence, communicate insights clearly and support improvement directions in high-constraint environments.

Relevant for \textbf{system performance analysis}, \textbf{software/hardware interaction diagnostics}, \textbf{advanced modelling}, \textbf{statistical analysis}, \textbf{simulation workflows}, \textbf{technical studies}, \textbf{validation}, \textbf{Unix-based engineering} and \textbf{collaborative multi-site environments}.

\cvsection{Fit for Arm System Performance}

\begin{itemize}
  \item Strong experience conducting deep technical studies on complex systems, performance limitations, uncertainty sources and degraded behaviours.
  \item Solid background in statistics, Bayesian inference, signal processing, modelling, simulation and sensitivity analysis.
  \item Experience identifying bottlenecks, deviations, numerical behaviours, weak points and possible improvement or mitigation directions.
  \item Hands-on experience with \textbf{Python}, \textbf{C/C++ foundations}, \textbf{Unix/Linux}, Bash scripting and Git/GitLab.
  \item Used to analysing system-level behaviours at the interface between software, data, modelling assumptions and operational constraints.
  \item Comfortable producing structured technical evidence, documenting assumptions and communicating results to multidisciplinary teams.
  \item Familiar with high-performance and high-constraint technical contexts through space mission analysis, precision instrumentation and detector performance studies.
  \item Motivated to extend existing system-performance expertise toward CPU, interconnect and cloud/server SoC environments.
\end{itemize}

\cvsection{Team Management}

\begin{itemize}
  \item Supervision of \textbf{3 Master's thesis interns (M2)} during engineering, modelling and software-development phases.
  \item Coordination of technical activities involving \textbf{research engineers and technicians} on a \textbf{data processing and validation pipeline} for the \textbf{LISA space mission}.
  \item Experience in organizing tasks, supporting contributors, reviewing outputs and maintaining alignment across collaborative technical developments.
\end{itemize}

\cvsection{Languages}

\cvskill{English}{4}
\divider

\cvskill{Spanish}{2}
\divider

\medskip

\cvsection{Education}

\cvevent{PhD in Physics}{Université Côte d'Azur}{2018 -- 2021}{Nice, France}
\divider

\cvevent{Selective Graduate Program in Fundamental Physics}{Université Paris-Saclay}{2016 -- 2018}{Orsay, France}
\divider

\cvevent{MSc in Astronomy and Astrophysics}{Observatoire de Paris / Université Paris-Saclay}{2017 -- 2018}{Paris, France}
\divider

\cvevent{BSc in Fundamental Physics}{Université Paris-Saclay}{2015 -- 2016}{Orsay, France}

\end{paracol}

\cvsection{Projects}

\cvevent{LISA Consortium -- Software, Modelling \& System Performance Analysis}{ESA/NASA Space Mission -- Large-scale International Collaboration}{2018 -- now}{}

Contribution to software, analysis and validation workflows for the LISA space mission within an international engineering and scientific collaboration involving ESA, NASA and multiple research institutes.

\textbf{Role:} Contribution to software workflows, technical coordination, complex signal-processing activities, model integration, validation campaigns, performance assessment and collaborative engineering tasks.

\textbf{System performance relevance:} Work involving simulation, model comparison, weak-signal analysis, sensitivity studies, uncertainty analysis, bottleneck identification, validation evidence and performance limitation assessment.

\textbf{Project relevance:} Work at the interface between software development, scientific requirements, data processing, validation workflows, technical contributors and international coordination.

\textbf{Program scale:} Major international space mission with an estimated budget of approximately 2 billion EUR over 10 years.

\divider

\cvevent{Einstein Telescope and LIGO--Virgo--KAGRA Collaborations}{International Scientific and Engineering Collaborations}{2021 -- now}{}

Contribution to collaborative developments in data analysis, signal characterization, software methods and technical investigations within the LIGO--Virgo--KAGRA and Einstein Telescope collaborations.

\textbf{Role:} Development of analysis tools, contribution to technical activities and support to complex software workflows in high-standard collaborative environments.

\textbf{System impact:} Contributed to studies on environmental and instrumental noise sources, supporting the identification of performance limitations and design constraints for next-generation gravitational-wave observatories.

\textbf{Environment:} Large-scale international collaborations involving strong technical requirements, structured methods, technical review, versioned software workflows and coordinated contributions.

\cvsection{Strengths}

\vspace{-.3cm}

\cvsubsection{System performance \& technical analysis}

{\small
\cvtag{System Performance Analysis}
\cvtag{Software Performance Analysis}
\cvtag{Performance Modelling}
\cvtag{Bottleneck Identification}
\cvtag{Deep Technical Studies}
\cvtag{Statistical Modelling}
\cvtag{Bayesian Inference}
\cvtag{Simulation Workflows}
\newline
\cvtag{Uncertainty Analysis}
\cvtag{Sensitivity Analysis}
\cvtag{Signal Processing}
\cvtag{Noise Modelling}
\cvtag{Technical Diagnostics}
\cvtag{Anomaly Analysis}
}

\divider\smallskip
\vspace{-.3cm}

\cvsubsection{Systems, validation \& diagnostics}

{\small
\cvtag{Systems Engineering}
\cvtag{Complex Systems}
\cvtag{Mission-Level Performance}
\cvtag{Validation Workflows}
\cvtag{Verification \& Validation}
\cvtag{Software Integration}
\cvtag{Integration Testing}
\newline
\cvtag{Requirements Traceability}
\cvtag{Configuration Follow-Up}
\cvtag{Root Cause Analysis}
\cvtag{Technical Evidence}
\cvtag{Software / Hardware Interfaces}
\cvtag{Operational Diagnostics}
}

\divider\smallskip
\vspace{-.3cm}

\cvsubsection{Technical skills}

{\small
\cvtag{Python}
\cvtag{C / C++}
\cvtag{Unix / Linux}
\cvtag{Bash}
\cvtag{Shell Scripting}
\cvtag{Git / GitLab}
\cvtag{Source Control}
\cvtag{Docker}
\cvtag{SQL}
\cvtag{Pandas}
\cvtag{Matlab}
\cvtag{TypeScript}
\cvtag{Node.js}
\cvtag{REST API}
\cvtag{FTP}
\cvtag{Jira}
\cvtag{Confluence}
\cvtag{OpenSearch}
}

\divider\smallskip
\vspace{-.3cm}

\cvsubsection{Engineering environment}

{\small
\cvtag{Reproducible Workflows}
\cvtag{Structured Documentation}
\cvtag{Technical Reporting}
\cvtag{Multi-Site Collaboration}
\cvtag{Cross-Functional Coordination}
\cvtag{International Collaboration}
\cvtag{Technical English}
\newline
\cvtag{Space Systems}
\cvtag{Precision Instrumentation}
\cvtag{Industrial Software}
\cvtag{Operational Monitoring}
\cvtag{Scientific Computing}
}


\end{document}

